\section{Scaling and reblocking}\label{subsec:rg-scaling}

Throughout Sections~\ref{subsec:rg-scaling}--\ref{subsec:rg-activity-contraction},
both $\ell$ and $\ell/L$ are admissible and the hypotheses of
Theorem~\ref{thm:main} hold. In particular $\ell\ge13L^2M$.
We first record the exact scaling identities from
\cite[\S12.6]{QED3Manuscript}.

Write the three thresholds of $R_k$ as
\[
 \theta_1=e_k^{-3/4+\varepsilon},\qquad
 \theta_2=e_k^\varepsilon\theta_1,\qquad
 \theta_3=e_k^\varepsilon\theta_2,
\]
and the thresholds of $\mathfrak X_k$ as
$d_1=e_k^{-1+2\varepsilon}$ and $d_2=e_k^\varepsilon d_1$.
For a level $k+1$ bond function and gauge function define
\begin{equation}\label{eq:rg-connection-scaling}
 A_L(b)=L^{-1/2}A(b/L),\qquad
 \lambda_L(c)=L^{1/2}\lambda(c/L).
\end{equation}
Since $\eta_{k+1}=\eta_k/L$, the finite differences obey
$(d\lambda)_L=d\lambda_L$ and
\[
 \|A_L\|_\infty=L^{-1/2}\|A\|_\infty,\qquad
 \|\partial A_L\|_\infty=L^{-3/2}\|\partial A\|_\infty,\qquad
 \|\delta_\alpha\partial A_L\|_\infty
                \le L^{-3/2-\alpha}\|\delta_\alpha\partial A\|_\infty.
\]
The last inequality uses only the coarse pairs at distance
$\le1/L$, a subset of those in the coarse seminorm.
The charge relation $e_{k+1}=L^{1/2}e_k$ therefore gives
the following fine-scale bounds for the scaled complex
small field, and for each scaled local real representative:
\begin{equation}\label{eq:rg-scaled-field-thresholds}
 \begin{split}
 b_1&=L^{-7/8+\varepsilon/2}\theta_1,\\
 b_2&=L^{-15/8+\varepsilon}\theta_2,\\
 b_3&=L^{-15/8+3\varepsilon/2}\theta_3.
 \end{split}
\end{equation}
Here the additional factor $L^{-\alpha}$ in the third bound
has been discarded. For the external argument the bounds are
\begin{equation}\label{eq:rg-scaled-external-thresholds}
 \|\mathscr C_L\|_\infty<L^{-1+\varepsilon}d_1,\qquad
 \|\partial\mathscr C_L\|_\infty<L^{-2+3\varepsilon/2}d_2.
\end{equation}
The enlarged coarse cube used to define a local representative,
after multiplication by $L$, contains the enlarged fine cube.
Scaling its real gauge thus gives the corresponding fine
representative. All the ratios in
\eqref{eq:rg-scaled-field-thresholds} and
\eqref{eq:rg-scaled-external-thresholds} are at most $L^{-3/4}$,
by \eqref{eq:rg-contraction-parameters}. It follows that
\begin{equation}\label{eq:rg-scaled-domain-inclusion}
 \mathscr A_L\in L^{-3/4}\widetilde R_k
                     \subset\widetilde R_k^-,\qquad
 \mathscr C_L\in L^{-3/4}\mathfrak X_k.
\end{equation}
For the second inclusion in the first statement use
$L^{-3/4}\le1/2\le1-e_k^\varepsilon$.

The Grassmann scaling homomorphism is
\begin{equation}\label{eq:rg-leg-scaling}
 S_L\psi(x)=L^{-1}\psi'(x/L),\qquad
 S_L\chi(x,x')=L^{-1-\alpha}\chi'(x/L,x'/L),
\end{equation}
with the same powers for the barred generators.
Only the photon coordinate is divided by $L$; the chamber
label and spin index are kept. Ordinary indices are mapped
bijectively onto the coarse ordinary indices. Pair indices
are mapped bijectively onto the coarse pair indices with
$d_X\le M/L$, since $d_X$ is divided by $L$.
This includes a pair inside one cluster:
$d_X=\eta_k$ becomes $d_X=\eta_{k+1}$.
The antiperiodic signs agree under this bijection.

Thus distinct basis monomials have distinct images, up to
their ordering signs, and every monomial of bidegree $(n,m)$
is multiplied by $L^{-n-(1+\alpha)m}$. There is no cancellation
of coefficients. For arbitrary positive weights,
\begin{equation}\label{eq:rg-scaling-norm-identity}
 \|S_LF\|_{w_1,w_2}=\|F\|_{L^{-1}w_1,L^{-1-\alpha}w_2}.
\end{equation}
As $h_{k+1,1}/h_{k,1}=L^{-1/8}$ and
$h_{k+1,2}/h_{k,2}=L^{-1/8+\varepsilon/2}$, this gives
\begin{equation}\label{eq:rg-degree-scaling}
 \|(S_LF)^{[n,m]}\|_{h_{k+1}}
       =L^{-p(n,m)}\|F^{[n,m]}\|_{h_k},\qquad
 p(n,m)=n+(1+\alpha)m+\frac{n+m}{8}-\frac{m\varepsilon}{2}.
\end{equation}
The powers are identical for all admissible pair types.

In kernel notation the same calculation reads
\begin{equation}\label{eq:rg-scaled-kernel}
 (S_LF)_{nm}(\xi';\zeta')=
 \begin{cases}
 L^{3n+6m-n-(1+\alpha)m}F_{nm}(L\xi';L\zeta'),
                       &d_X(\zeta'_j)\le M/L\text{ for every }j,\\
 0,                    &\text{otherwise}.
 \end{cases}
\end{equation}
Indeed, the index weights satisfy
$\eta_k^3/48=L^3\eta_{k+1}^3/48$ and hence contribute
$L^{3n+6m}$ to the kernel. Antisymmetry and antiperiodicity
are preserved, so uniqueness of the kernel expansion gives
\eqref{eq:rg-scaled-kernel}.

Let $\bar X^L$ be the union of $LM$-cubes containing the cubes
of $X$. Define the reblocked family by
\begin{equation}\label{eq:rg-reblocking-definition}
 (\mathcal L F)(Y;\mathscr A,\mathscr C)
       =\sum_{X:\bar X^L=LY}S_L
                          F(X;\mathscr A_L,\mathscr C_L).
\end{equation}
The sum is finite on a fixed torus. It is local in $Y$ if
$F$ is local in each $X$. Ordinary legs over $X$ scale to
legs over $Y$; the six paths for a pair of the first kind
scale to paths in $Y$. For a boundary pair,
$x^{(3)}/L=(x/L)^{(3)}$, and an endpoint in $X$ scales to
an endpoint in $Y$. Finally, every fine bond meeting $X$
scales to a coarse bond meeting $Y$. Equality of both
coarse connection arguments on these bonds consequently
gives equality of the fine arguments read by each term.
This proves both the support and the coefficient locality
of \eqref{eq:rg-reblocking-definition}.
